% This is the ADASS_template.tex LaTeX file, 24 Oct 2025.
% It is based on the ASP general author template file, but modified to reflect the specific
% requirements of the ADASS proceedings.

% This is the aspauthor.tex LaTeX file
% Copyright 2021, Astronomical Society of the Pacific Conference Series
% Revision:  29 November 2021

% To compile, at the command line positioned at this folder, type:
% pdflatex ADASS_template
% pdflatex ADASS_template
% This will create a file called ADASS_template.pdf

\documentclass[11pt,twoside]{article}

% Do NOT use ANY packages other than asp2023. 
\usepackage{asp2023}

\aspSuppressVolSlug
\resetcounters

% References must all use BibTeX entries in a .bibfile.
% References must be cited in the text using \citet{} or \citep{}.
% Do not use \cite{}.
% See ManuscriptInstructions.pdf for more details
\bibliographystyle{asp2021}

% The ``markboth'' line sets up the running heads for the paper.
% 1 author: "Surname"
% 2 authors: "Surname1 and Surname2"
% 3 authors: "Surname1, Surname2, and Surname3"
% >3 authors: "Surname1 et al."
% Replace ``Short Title'' with the actual paper title, shortened if necessary.
% Use mixed case type for the shortened title
% Ensure shortened title does not cause an overfull hbox LaTeX error
% See ASPmanual2010.pdf 2.1.4  and ManuscriptInstructions.pdf for more details
\markboth{Danehkar}{xstar-tools: AI-Assisted Modernization of XSTAR}

\begin{document}

\title{xstar-tools: AI-Assisted Modernization of the XSTAR Photoionization Code}

% Note the position of the comma between the author name and the 
% affiliation number.
% Authors surnames should come after first names or initials, eg John Smith, or J. Smith.
% Author names should be separated by commas.
% The final author should be preceded by "and".
% Affiliations should not be repeated across multiple \affil commands. If several
% authors share an affiliation this should be in a single \affil which can then
% be referenced for several author names. If only one affiliation, no footnotes are needed.
% See ManuscriptInstructions.pdf and ASP's manual2010.pdf 3.1.4 for more details
\author{Ashkbiz~Danehkar$^1$}
\affil{$^1$Universities Space Research Association, Huntsville, AL, 35805, USA}

% This section is for ADS Processing.  There must be one line per author. paperauthor has 9 arguments.
\paperauthor{Ashkbiz Danehkar}{adanehkar@usra.edu}{0000-0003-4552-5997}{Universities Space Research Association}{Science and Technology Institute}{Huntsville}{AL}{35805}{USA}

% There should be one \aindex line (commented out) for each author. These are used to
% build up the author index for the Proceedings. The surname must come first, followed by
% initials. Note the use of ~ before each initial to control spacing.
% The \author entries (see above) have surname last. These \aindex entries have
% surname first.
% The Aindex.py command willl create them for you after you have constructed the \author
% The first entry should be the first author, for bold-facing the author index page-reference

%\aindex{FistAuthor1,~S.~A.}
%\aindex{Author2,~S.~B.}
%\aindex{Author3,~S.}


\begin{abstract}
We present {\tt xstar-tools}, an open-source modernization of the \texttt{Fortran} \textsc{xstar} photoionization code developed through iterative, AI-assisted software development. The suite provides an independent C{}\texttt{++} engine, Python interfaces, and native workflows for generating \textsc{xspec}-compatible spectral tables. Large language models assisted translation, discrepancy analysis, optimization, testing, and documentation under human supervision. Scientific qualification against \texttt{Fortran} \textsc{xstar} examined atomic populations, heating and cooling, convergence, radiative transfer, and spectra. Numerical acceptance criteria were complemented by regression checks against accepted implementation baselines. For a broad multi-element model, AI-assisted optimization substantially reduced runtime and peak memory relative to earlier C{}\texttt{++}, although performance remains model-dependent. Serial, local multiprocess, and optional MPI grid execution support scientific workflows across supported platforms. We discuss reference implementations, intermediate diagnostics, and human review as foundations for trustworthy modernization, with future development extending the framework to post-processing hydrodynamic simulations.
\end{abstract}

% These lines show examples of subject index entries. At this stage these have to commented
% out, and need to be on separate lines. Eventually, they will be automatically uncommented
% and used to generate entries in the Subject Index at the end of the Proceedings volume.
% Don't leave these in! - replace them with ones relevant to your paper.

% These lines show examples of ASCL index entries. At this stage these have to commented
% out, and need to be on separate lines. Eventually, they will be automatically uncommented
% and used to generate entries in the ASCL Index at the end of the Proceedings volume.
% The ascl.py command will scan your paper on possible code names.
% Don't leave these in! - replace them with ones relevant to your paper.
%\ooindex{FOOBAR, ascl:1101.010}


\section{Scientific Background and Modernization Goals}
\textsc{xstar} is an established \texttt{Fortran} code for calculating the physical conditions and emitted and transmitted spectra of photoionized gas. Given an incident radiation field, abundances, and density, it calculates ionization balance, level populations, opacities, and line and continuum emission. It combines radiative and collisional processes with radiative transfer, solving thermal balance or using a prescribed temperature \citep{Kallman2001,Kallman2004}. The underlying atomic database includes the transition probabilities, cross sections, and rate coefficients for these calculations \citep{Bautista2001}.

We present {\tt xstar-tools}, an open-source modernization developed through iterative, AI-assisted software development under human supervision. The project progressed from a pure Python implementation to an independent C{}\texttt{++} engine with Python interfaces and native workflows for planning model grids, executing calculations, and assembling \textsc{xspec}-compatible spectral tables. The aim is to improve portability, maintainability, and integration while retaining the established \textsc{xstar} physical framework.

Changing implementation language and data representation introduces challenges in indexing, numerical evaluation, memory ownership, and output organization. We therefore evaluate scientific agreement, runtime and memory use, and output compatibility through separate qualification checks. \texttt{Fortran} \textsc{xstar} provides the scientific reference, while accepted implementation baselines guard against unintended changes during subsequent development.

\section{AI-Assisted Software Development and Scientific Qualification}
Large language models assisted translation, discrepancy analysis, optimization, benchmarks, and documentation. Human reviewers checked numerical semantics and physical assumptions. Development proceeded through bounded revisions: identify a discrepancy or expensive operation, inspect relevant source behavior, implement a candidate, and evaluate it before promotion. This workflow gave AI-generated changes a concrete reference and a reviewable acceptance path (Figure~\ref{fig:workflow}).

\articlefigure[width=.8\textwidth]{fig1_xstar-tools.eps}{fig:workflow}{AI-assisted development workflow combining human review, scientific comparisons with \texttt{Fortran} \textsc{xstar}, regression testing, and performance checks.}

\texttt{Fortran} \textsc{xstar} served as the authoritative scientific reference; the historical optimization campaign used version 2.59g. 
Qualification examined atomic level populations, heating and cooling rates, convergence trajectories, radiative transfer, emergent spectra, and output products. Intermediate-state comparisons were especially useful when a final spectrum alone could not identify where a discrepancy arose. They allowed investigations to follow the calculation from atomic-record interpretation through population solutions to spectral publication.

Two comparison modes served different purposes. Numerical acceptance criteria assessed agreement between C{}\texttt{++} and \texttt{Fortran}. Exact regression checks against accepted C{}\texttt{++} baselines detected unintended changes within the implementation lineage. The latter can include selected FITS payloads and normalized STEP logs, with explicitly identified metadata treated separately. Exact regression agreement must not be interpreted as a general claim of bit-for-bit agreement between C{}\texttt{++} and \texttt{Fortran}. Cross-language tolerances must be attached to the quantity and test being evaluated.

Iteration counts and thermal balance diagnostics help compare convergence trajectories. Small residuals in a reference log describe that calculation; they do not measure cross-language discrepancies. Qualification combines these diagnostics with population, radiation-field, and spectral comparisons.

\section{Memory and Runtime Optimization}
Profiling and ownership analysis guided the performance campaign. Important changes included compact representations of atomic-record products, storage of only required bound-free data spans, shorter lifetimes for transient state, and removal of duplicate payloads and intermediate copies. Runtime work also reused prepared rate kernels and continuum workspaces and removed repeated administrative calculations from frequently executed paths. These changes targeted implementation overhead while preserving the scientific calculation and its required ordering.

Table~\ref{tab:laptop} summarizes the historical same-machine comparison for the broad multi-element test identified in the project records as \texttt{multi\_element\_xi1\_ne1e12} ($\log \xi=1$ and $n_{\rm H}=10^{12}$ cm$^{-3}$). The earlier C{}\texttt{++} baseline required approximately two hours, whereas the accepted optimized base required approximately 12 minutes. The summarized peak-memory values fell from about 11--12~GB to about 3.7~GB. Thus, the campaign produced an approximately tenfold runtime improvement and reduced peak memory by roughly two-thirds relative to the earlier C{}\texttt{++} implementation.

\begin{table}[ht]
\caption{Historical C{}\texttt{++} optimization comparison on the same machine.}
\label{tab:laptop}
\begin{center}\small
\begin{tabular}{lccc}
\tableline
Metric & Earlier C{}\texttt{++} & Optimized C{}\texttt{++} & Improvement\\
\tableline
Wall time & $\sim120$ min & $\sim12$ min & $\sim10\times$ faster\\
Peak memory & $\sim11$--12 GB & $\sim3.7$ GB & $\sim2/3$ lower\\
\tableline
\end{tabular}
%\\
%\smallskip
\end{center}
\noindent{\footnotesize \textbf{Note.} Earlier: 0.6.82.33.8; Optimized: 0.6.82.40.2.46.1. Values are approximate historical summaries on an HP laptop (Intel Core Ultra 5 135U, 32 GB Memory); memory units are retained as reported.}
\end{table}

A separate multi-element comparison on a workstation is given in Table~\ref{tab:workstation}. The optimized C{}\texttt{++} model completed in 661~s, compared with 941~s for \texttt{Fortran}: a 29.8\% reduction in wall time, or a speedup of 1.42. Peak RSS remained higher for C{}\texttt{++}. %, at 3.52~GB versus 2.40~GB. 
%These RSS values were converted consistently from the archived Linux readings of 3,691,864 and 2,513,224~KB, respectively. 
The optimization reduced the C{}\texttt{++} memory deficit substantially, but did not eliminate it.

\begin{table}[ht]
\caption{Historical multi-element comparison on the same machine.}
\label{tab:workstation}
\begin{center}\small
\begin{tabular}{lcc}
\tableline
Metric & \texttt{Fortran} \textsc{xstar} & Optimized C{}\texttt{++}\\
\tableline
Wall time & 941 s & 661 s\\
Peak RSS & 2.57 GB & 3.78 GB\\
\tableline
\end{tabular}
\end{center}
\noindent{\footnotesize \textbf{Note.} Values are for runs on an ASUS workstation (AMD Ryzen 9 5950X, 64 GB Memory). 
%Memory is expressed in decimal GB ($10^9$ bytes).
}
\end{table}

These results characterize the recorded model and hosts, not a universal speed advantage. Low-ionization cases remained slower than \texttt{Fortran}, and runtime and memory depend on the active elements and model conditions. Different machine records are separate comparisons and should not be combined into a single timing comparison. They are historical measurements, not repeated statistical benchmarks. % with uncertainty estimates. 
A comprehensive benchmark release should include complete input files, atomic-data versions, compiler flags, hardware specifications, raw timing logs, and repeated measurements.

\section{Native Model-Grid Workflows and Deployment}
The native workflow connects parameter-grid planning, independent \textsc{xstar} model execution, and \textsc{xspec}-compatible table assembly. It supports serial execution, independent local worker processes, and optional MPI execution on high-performance computing clusters, following the parallelization approach of earlier work \citep{Danehkar2018}. Parallelism is applied across model jobs; it does not imply that the numerical calculation inside each model has been parallelized. Isolated job outputs also make individual failures easier to identify and investigate.

Results are collected in canonical grid order, so completion order does not determine placement in the spectral table. Model FITS products and ordered STEP-log concatenation are handled separately. Table assembly combines scientific arrays; log assembly preserves a traceable grid record. Both transformations are qualified separately from the model engine.

Native builds have been qualified on Linux, macOS, and Windows. Distribution uses PyPI and conda-forge on their supported platforms; Windows native support uses the MinGW/PyPI path. Conda distribution and MPI are not supported on Windows, and MSVC support is not planned. Process concurrency must account for the several-gigabyte memory footprint of broad models, rather than defaulting unconditionally to the number of available CPU cores. Source, installation instructions, and documentation are available from the project repository.\footnote{\url{https://github.com/xstar-tools/xstar-tools}}

\section{Lessons and Future Work}
The main lesson is that AI assistance is most useful when proposals are constrained by inspectable references and explicit acceptance criteria. Compilation and plausible outputs are necessary but insufficient evidence for scientific software. Intermediate diagnostics help turn a mismatch into a localized investigation; retained baselines make subsequent changes reviewable. Human supervision remains responsible for physical interpretation, acceptance thresholds, and decisions about whether an optimization is scientifically admissible.

Compact storage and improved ownership can reduce memory traffic and runtime without changing the physical framework. Separate scientific and performance checks make model-dependent tradeoffs visible.

Future work will reduce intermediate file I/O while retaining the file-backed workflow as a compatibility reference. We also plan to extend the framework toward post-processing hydrodynamic and magnetohydrodynamic simulations, using simulated gas conditions to calculate synthetic spectra. This extension will require explicit assumptions about the incident radiation, geometry, and applicability of equilibrium calculations. It remains planned development rather than a capability established by the benchmarks presented here.

\bibliography{P304}  % For BibTex

% if we have space left, we might add a conference photograph here. Leave commented for now.
% \bookpartphoto[width=1.0\textwidth]{foobar.eps}{FooBar Photo (Photo: Any Photographer)}

\end{document}
